\begin{table}[htbp]
  \centering
  \small
  \begin{tabular}{l|ccc}
    \toprule
    \textbf{Belief Conditioning} & \textbf{Gemma 2 27B} & \textbf{Llama 3.1 70B} & \textbf{Llama 3.1 8B} \\
    \midrule
    Unconditioned & 0.40 & 0.50 & 0.40 \\
    Self-Conditioned & 0.00 & 0.80 & 1.00 \\
    \cmidrule{1-4}
    Weak Perturbation ($\rho=0.80$) & 0.08 & 0.30 & -0.14 \\
    Strong Perturbation ($\rho=0.20$) & 0.14 & 0.20 & 0.40 \\
    \bottomrule
  \end{tabular}
  \caption{
    Belief conditioning effectiveness across models and methods. Values are Spearman correlations ($\rho$; higher is better) between the imposed prior and simulated behavior. \textbf{Unconditioned} reports baseline consistency. \textbf{Self-conditioned} supplies each model's own elicited beliefs as context. \textbf{Weak} and \textbf{strong perturbation} conditions apply modified priors constructed to have $\rho=0.80$ or $\rho=0.20$ correlation, respectively, with the original elicited beliefs. Self-conditioning shows highly model-dependent effects; even weakly perturbed priors can substantially disrupt belief--behavior alignment.
  }
  \label{tab:belief-conditioning}
\end{table}
